\section{Introduction}

The design of optical systems has long relied on computational tools for evaluating and refining candidate lens assemblies. As outlined in \cite{Smith2008ModernOpticalEngineering}, lens design workflows are inherently iterative, involving repeated analysis of ray behavior, aberrations, and other system performance-and-manufacturability metrics, followed by numerical optimization of design parameters. Modern optical design environments like \cite{codev} and \cite{zemax_opticstudio} integrate high fidelity ray tracing engines with local and global numerical optimization algorithms, enabling refinement of optical system parameters with respect to a set of chosen performance objectives.

Despite their capabilities, the effectiveness of such tools remains strongly dependent on the quality of starting design configuration and its parameters. Optical design performance space is characterized by highly non-convex landscapes with numerous poor local minima, making optimization tools sensitive to initialization \cite{Turnhout:09}. Furthermore, most commercial tools assume fixed system topology during optimization, leaving exploration across different number of elements or alternative lens and stop surface configurations largely as a manual task that requires substantial domain expertise. As a result early-stage design remains a major bottleneck, with design engineers spending significant effort manually searching and adapting existing patent and internal database designs to obtain viable starting point before refining them for current system requirements.

Recent progress in deep learning has enabled data-driven exploration of complex engineering design spaces \cite{regenwetter2022deepgenerativemodelsengineering}. Diffusion-based generative models, in particular, have been applied to a range of design synthesis and inverse design problems. For example, latent diffusion models have been used for CAD sequence reconstruction in \cite{alam2025gencadimageconditionedcomputeraideddesign}.
Usage of conditional diffusion models for inverse design of blended wing bodies is explored in \cite{blendednet_pp_Sung_2025} and their application in ship hull design exploration is studied in \cite{bagazinski2023shipgendiffusionmodelparametric}

In optical system design, deep learning has been explored for a range of tasks, including optimization acceleration, inverse design, and starting point generation. For example, surrogate models have been applied to gradient-based optimization of thin film designs in \cite{hegde2019accelerating}. Starting point generation for free-form reflective triplet systems using deep neural networks is explored in \cite{Yang:19}. More closely related to lens assembly design, recurrent neural network architectures have been used for design space exploration in \cite{Cote:21}. However, this approach relies on relatively simple deep learning models and is limited in several respects, including restriction to unit effective focal length, manual specification of glass-air topology and stop surface location, and training on a small corpus of optical designs and operating conditions (approximately 150 systems).

These prior efforts suggest that generative modeling has the potential to support optical system design, but also highlight challenges that are specific to this domain. Optical systems are naturally represented as variable-length sequences of elements, and many combinations of parameters can lead to non-physically realizable designs or unrealistic material choices. As a result, careful consideration must be given to how optical designs are represented for learning and how generated prescriptions are post-processed and evaluated before they are treated as usable optical designs.

Building on these observations, this work makes four contributions toward generative optical lens design. First, we compile a large-scale dataset of patented and synthetically expanded optical systems with associated performance metrics. Second, we introduce a parametric representation of lens assemblies designed to improve generative modeling over direct raw prescription tables. Third, we benchmark off-the-shelf unconditional generative models for lens prescription generation. Fourth, we incorporate a ray-tracing-based post-processing pipeline for improving the physical feasibility, simulation readiness, and optical performance of generated samples.

The remainder of this paper is organized as follows. We first review related work in optical design automation and generative modeling for engineering design. We then describe the patent-derived and synthetic lens datasets, followed by the proposed representation and post-processing pipeline. Next, we present the benchmark models, evaluation framework, and experimental results. We conclude with discussion of limitations and directions for future work.
